\documentclass[11pt]{article}
\usepackage[a4paper, margin=1in]{geometry}
\usepackage[T1]{fontenc}
\usepackage[utf8]{inputenc}
\usepackage{amsmath, amssymb}
\usepackage{booktabs}
\usepackage{hyperref}
\usepackage{xcolor}
\hypersetup{colorlinks=true, linkcolor=blue, urlcolor=blue, citecolor=blue}
\usepackage{authblk}
\usepackage{microtype}
\usepackage{enumitem}
\setlist[itemize]{leftmargin=1.4em}

\title{\textbf{RealQM vs.\ DFT for Protein Folding}\\[0.3em]
\large Why density-functional theory cannot dynamically simulate folding, and RealQM can}
\author[1]{Claes Johnson}
\affil[1]{KTH Royal Institute of Technology, Stockholm, Sweden \\ \texttt{claesjohnson@gmail.com}}
\date{\today \\[0.4em]
\small\textit{Developed in collaboration with Claude (Anthropic).}}

\begin{document}
\maketitle

\begin{abstract}
Protein folding is a \textbf{dynamical} process --- a chain moving from an extended state to its native fold
over microseconds, in a system of $10^3$--$10^5$ atoms. We argue that \textbf{density-functional theory (DFT)
cannot follow this trajectory}: its cost scales as $N^3$ in the electron number and it is locked to
femtosecond steps, putting a microsecond fold of a solvated protein at ${\sim}10^{13}$ core-hours
(${\sim}$millennia) --- $10^9$--$10^{12}$ beyond feasibility. \textbf{RealQM can}: its cost is set by a
\emph{fixed spatial grid} ($200^3$--$1000^3$ cells), \emph{independent of the atom count} (a $200^3$ grid
handles ${\sim}1{,}000$ atoms, a $1000^3$ grid ${\sim}100{,}000$), and it runs large-step (Brownian) dynamics
on the GPU --- with folding trajectories already demonstrated (chignolin, crambin, GB1, solvated $\alpha$-helix,
coiled coil). Biases may be used to guide the fold; that helps RealQM and does nothing for DFT, whose
disqualifying cost is \emph{per step}.
\end{abstract}

\section{The problem: folding is dynamics, at scale}
Folding is not a single minimisation but a \textbf{time-evolution} through a funnel of conformations.
Timescale: microseconds (small fast-folders) to milliseconds. Size: ${\sim}10^3$ atoms bare, but
$10^4$--$10^5$ atoms (${\sim}10^5$ electrons) with explicit solvent. What is required is to follow the
\emph{whole trajectory}, not to score a static structure --- so any quantum method must run \textbf{many time
steps on a large system}. This is where DFT and RealQM diverge.

\section{Why DFT cannot}
\begin{itemize}
\item \textbf{$O(N^3)$ in electrons.} Each self-consistent solve scales as the cube of the electron number
(orthogonalisation / diagonalisation); a solvated protein (${\sim}10^5$ e) is ${\sim}10^6\times$ a
$1{,}000$-atom benchmark.
\item \textbf{Femtosecond steps.} Ab-initio MD is locked to ${\sim}1$ fs by nuclear vibrations, so
$1\,\mu\text{s}=10^9$ steps, each a full electronic solve.
\item \textbf{Bottom line.} For ${\sim}5\,\mu$s of a solvated protein: $5\times10^9$ steps $\times\ {\sim}10^6$
core-seconds/step $\approx \mathbf{10^{13}}$ \textbf{core-hours} $\approx$ \textbf{${\sim}1{,}000+$ years on a
million-core machine} --- for one trajectory.
\end{itemize}
In practice DFT-AIMD reaches ${\sim}100$ ps for ${\sim}1{,}000$ atoms. Folding needs ${\sim}10^6\times$ longer
\emph{and} ${\sim}30\times$ larger ($\times N^3=10^4\times$ per step) --- a combined
$\mathbf{{\sim}10^{9}\text{--}10^{12}}$ beyond feasibility, a wall rather than a hardware wait. \textbf{Biases
do not help DFT}: restraints cut how many steps you must \emph{sample}, but the disqualifying cost is the
\emph{per-step} $N^3$ electronic solve, paid every step.

\section{Why RealQM can}
RealQM is quantum mechanics as real-space continuum mechanics on a grid, with a different cost structure:
\begin{itemize}
\item \textbf{Cost set by the grid, not by $N$.} All atoms share \emph{one fixed spatial grid}; per-step cost
scales with the number of grid cells (multigrid), \emph{independent of atom count}. A $200^3$ grid handles
${\sim}1{,}000$ atoms; a $1000^3$ grid ${\sim}100{,}000$ atoms --- at the same per-cell cost. Adding atoms does
not blow the cost up as $N^3$ does.
\item \textbf{Large-step dynamics.} Overdamped / Brownian dynamics replaces fs-locked Newtonian AIMD, reaching
folding timescales in orders of magnitude fewer steps.
\item \textbf{GPU-native.} The real-space grid maps directly onto GPU hardware.
\item \textbf{Demonstrated.} Folding trajectories already run --- chignolin, crambin, GB1, a solvated
$\alpha$-helix, a coiled coil.
\item \textbf{Biases permitted.} As in classical MD (native-contact biasing, restraints, enhanced sampling),
biases may guide the fold; this helps RealQM and, again, does nothing for DFT.
\end{itemize}

\section{Side by side}
\begin{center}
\begin{tabular}{lll}
\toprule
 & \textbf{DFT (ab-initio MD)} & \textbf{RealQM} \\
\midrule
cost per step   & $O(N^3)$ in electrons        & $O(\text{grid})$ --- fixed by resolution \\
capacity        & ${\sim}1{,}000$ atoms        & $200^3\!:1{,}000$; $1000^3\!:100{,}000$ atoms \\
time stepping   & fs-locked, $10^9$/$\mu$s     & Brownian / large-step \\
reach today     & ${\sim}100$ ps / $1{,}000$ atoms & full folding trajectories \\
${\sim}5\,\mu$s fold & ${\sim}10^{13}$ core-hr (millennia) & days--weeks on GPUs \\
biases          & don't help                  & permitted, guide the fold \\
verdict         & \textbf{impossible}         & \textbf{feasible} \\
\bottomrule
\end{tabular}
\end{center}
The gap between the verdicts is ${\sim}10^{6}$--$10^{10}$ in compute: the difference between ``not in any
foreseeable machine'' and ``runs on a GPU cluster.''

\section{Honest status}
Beyond the capability gap, RealQM has a real advantage on the \textbf{physics of the dynamics} that should be
stated plainly: it computes the \textbf{forces on the atoms directly from the electron charge potentials} ---
the actual Coulomb forces that drive the motion --- rather than as \textbf{gradients of an energy surface}, as
in the standard energy-first approach. Dynamics \emph{is} force; RealQM supplies force physically and directly.
On this basis RealQM is \textbf{stronger, not weaker} --- the force computation is a strength, not a weak point.

The one genuinely open question is quantitative \textbf{fidelity of the folded state} --- whether the physical
force field carries the chain to the correct native structure across the whole funnel, and how far biases
(native-contact restraints, enhanced sampling, standard tools in classical MD) are used. That is ordinary
validation, not a defect of the method: the physical, charge-based forces are a reason for confidence, and
biases help RealQM while doing nothing for DFT, which remains disqualified by cost.

\end{document}
